\documentclass[11pt]{article}

\usepackage[margin=1in]{geometry}
\usepackage{amsmath, amssymb}
\usepackage{xcolor}
\usepackage{listings}
\usepackage{hyperref}
\usepackage{array}
\usepackage{booktabs}

\hypersetup{
    colorlinks=true,
    linkcolor=blue!60!black,
    urlcolor=blue!60!black,
    citecolor=blue!60!black
}

\lstdefinestyle{bugpython}{
  language=Python,
  basicstyle=\ttfamily\small,
  keywordstyle=\color{blue!60!black},
  stringstyle=\color{green!40!black},
  commentstyle=\color{gray!70!black},
  showstringspaces=false,
  breaklines=true,
  frame=single,
  rulecolor=\color{gray!30},
  columns=fullflexible,
  keepspaces=true
}

\title{SymPy Bug Report}
\author{}
\date{}

\begin{document}

\maketitle

\section{Bug 19: \texttt{rsolve} returns a solution that violates supplied initial conditions}

\subsection{Minimal Reproducer}

\medskip

\begin{lstlisting}[style=bugpython]
from sympy import Function, rsolve, symbols

n = symbols("n", integer=True)
y = Function("y")

rec = y(n + 2) - (n + 2)*y(n + 1) + y(n)
sol = rsolve(rec, y(n), {y(0): 1, y(1): 1})

print("rsolve output:", sol)
print("returned y(0):", sol.subs(n, 0))
print("returned y(1):", sol.subs(n, 1))
\end{lstlisting}

\subsection{Observed Behavior}

\medskip

\begin{lstlisting}[style=bugpython]
rsolve output: 0
returned y(0): 0
returned y(1): 0
manual first terms: [1, 1, 1, 2, 7, 33]
\end{lstlisting}

SymPy returns the zero sequence.  Direct substitution into that returned expression gives
\(y(0)=0\) and \(y(1)=0\), even though the call explicitly supplied
\(y(0)=1\) and \(y(1)=1\).  The result is therefore not merely an incomplete
closed form or a different representation of the intended sequence; it violates
the constraints passed to the public API.

\subsection{Expected Behavior}

The solver should return a sequence satisfying the recurrence
\[
    y(n+2) - (n+2)y(n+1) + y(n) = 0
\]
and the initial conditions \(y(0)=1\), \(y(1)=1\), or else report that it cannot
solve the recurrence with those initial conditions.  It should never return the
zero sequence for this call.

\subsection{Mathematical Explanation}

The recurrence is equivalent to
\[
    y(n+2) = (n+2)y(n+1) - y(n).
\]
Starting from the supplied initial values gives
\[
\begin{aligned}
    y(0)&=1,\\
    y(1)&=1,\\
    y(2)&=2y(1)-y(0)=1,\\
    y(3)&=3y(2)-y(1)=2,\\
    y(4)&=4y(3)-y(2)=7,\\
    y(5)&=5y(4)-y(3)=33.
\end{aligned}
\]
Thus a valid solution with these initial conditions begins
\[
    1,\ 1,\ 1,\ 2,\ 7,\ 33,\ldots .
\]
The zero sequence does satisfy the homogeneous recurrence, but it satisfies
\(y(0)=0\) and \(y(1)=0\), not the supplied nonzero initial conditions.  Hence
SymPy's answer is a solution of the recurrence alone, not a solution of the
initial-value problem requested by the user.

\subsection{Source-Level Diagnosis}

The diagnosis identifies the relevant implementation as
\nolinkurl{sympy/solvers/recurr.py}, in the public \texttt{rsolve} function
around lines 816--841.  The traced call path is
\texttt{sympy.rsolve(...)} to coefficient extraction, then to
\texttt{rsolve\_hyper(coeffs, 0, n, symbols=True)} near line 806, followed by
the initial-condition handling in \texttt{rsolve}.  For the representative
recurrence, the traced internal call
\texttt{rsolve\_hyper([1, -(n + 2), 1], 0, n, symbols=True)} returns the
constant-free candidate \texttt{0} with an empty list of arbitrary constants.

The problematic rule is that initial conditions are only applied when the
candidate expression contains symbols for arbitrary constants:

\begin{lstlisting}[style=bugpython]
if symbols and init is not None:
    ...
\end{lstlisting}

When \texttt{symbols} is empty, this block is skipped.  In the failing case
\texttt{solution} is already the constant-free expression \texttt{0}, so
\texttt{rsolve} never forms or checks the equations
\texttt{solution.subs(n, 0) - 1} and \texttt{solution.subs(n, 1) - 1}.  It then
returns the candidate unchanged.  This source-level behavior causes the
mathematical failure because a constant-free candidate can still contradict
explicit initial conditions.  The diagnosis supports a plausible fix direction:
when initial conditions are provided, \texttt{rsolve} should validate
constant-free solutions as well, returning \texttt{None} or the same unsolved
initial-condition failure used elsewhere when any supplied condition is false.

\subsection{Additional Failing Instances}

The independent verification checks the family
\[
    y(n+2) - (n+a)y(n+1) + y(n) = 0,\qquad y(0)=1,\quad y(1)=1,
\]
for \(a=2,3,4,5,6,7\).  In every listed case SymPy returns the zero sequence,
so the returned values at \(n=0\) and \(n=1\) are both \(0\).  The expected
behavior is the same across the family: the returned expression must satisfy
the two supplied initial values, or the solver should report that it cannot
produce such a solution.

\begin{center}
\small
\renewcommand{\arraystretch}{1.3}
\begin{tabular}{@{}>{\raggedright\arraybackslash}p{0.44\linewidth}>{\raggedright\arraybackslash}p{0.23\linewidth}>{\raggedright\arraybackslash}p{0.23\linewidth}@{}}
\toprule
\textbf{Input / Expression} & \textbf{SymPy behavior} & \textbf{Expected behavior} \\
\midrule
\(a=2:\ y(n+2)-(n+2)y(n+1)+y(n)=0\), \(y(0)=y(1)=1\) & returns \(0\); returned \(y(0)=0\), \(y(1)=0\) & satisfy \(y(0)=y(1)=1\); forward terms \(1,1,1,2,7,33\) \\
\(a=3:\ y(n+2)-(n+3)y(n+1)+y(n)=0\), \(y(0)=y(1)=1\) & returns \(0\); returned \(y(0)=0\), \(y(1)=0\) & satisfy \(y(0)=y(1)=1\); forward terms \(1,1,2,7,33,191\) \\
\(a=4:\ y(n+2)-(n+4)y(n+1)+y(n)=0\), \(y(0)=y(1)=1\) & returns \(0\); returned \(y(0)=0\), \(y(1)=0\) & satisfy \(y(0)=y(1)=1\); forward terms \(1,1,3,14,81,553\) \\
\(a=5:\ y(n+2)-(n+5)y(n+1)+y(n)=0\), \(y(0)=y(1)=1\) & returns \(0\); returned \(y(0)=0\), \(y(1)=0\) & satisfy \(y(0)=y(1)=1\); forward terms \(1,1,4,23,157,1233\) \\
\(a=6:\ y(n+2)-(n+6)y(n+1)+y(n)=0\), \(y(0)=y(1)=1\) & returns \(0\); returned \(y(0)=0\), \(y(1)=0\) & satisfy \(y(0)=y(1)=1\); forward terms \(1,1,5,34,267,2369\) \\
\(a=7:\ y(n+2)-(n+7)y(n+1)+y(n)=0\), \(y(0)=y(1)=1\) & returns \(0\); returned \(y(0)=0\), \(y(1)=0\) & satisfy \(y(0)=y(1)=1\); forward terms \(1,1,6,47,417,4123\) \\
\bottomrule
\end{tabular}
\end{center}

\subsection{Related Issues Assessment}

A re-audit of the deduplication check identifies open public issues that report
the same underlying \texttt{rsolve\_hyper}-returns-\(0\) failure family, in which
\texttt{rsolve} silently returns the zero sequence for variable-coefficient /
higher-order recurrences:
\begin{itemize}
  \item \href{https://github.com/sympy/sympy/issues/27902}{sympy/sympy\#27902},
  ``rsolve: Gives 0 for higher order recurrences'' (open) --- the most direct
  match: its example \texttt{rsolve(y(n+2) - n*y(n+1) - y(n), y(n))} returns
  \(0\), the same second-order variable-coefficient symptom.
  \item \href{https://github.com/sympy/sympy/issues/17982}{sympy/sympy\#17982},
  ``Wrong result from rsolve'' (open).
  \item \href{https://github.com/sympy/sympy/issues/11063}{sympy/sympy\#11063},
  ``Some wrong answers from rsolve'' (open).
\end{itemize}
The deduplication verdict is therefore \emph{family known, specific instance
new}: the broad \texttt{rsolve}-returns-\(0\) failure is publicly reported,
while bug-019's specific lens --- the constant-free \(0\) candidate causing the
\texttt{if symbols and init is not None} initial-condition block to be skipped,
so supplied conditions are never validated --- is the new specific instance
within that known family.  The issue remains individually actionable because it
has a compact reproducer, a localized source-level mechanism, and a direct
regression test for the supported initial-condition API.

\subsection{Proposed SymPy Regression Test}

\medskip

\begin{lstlisting}[style=bugpython]
import pytest
from sympy import Function, S, rsolve, symbols


CASES = [2, 3, 4, 5, 6, 7]


@pytest.mark.parametrize("a", CASES)
def test_rsolve_constant_free_solution_satisfies_initial_conditions(a):
    n = symbols("n", integer=True)
    y = Function("y")
    rec = y(n + 2) - (n + a)*y(n + 1) + y(n)
    sol = rsolve(rec, y(n), {y(0): S.One, y(1): S.One})

    assert sol is not None
    assert sol.subs(n, 0) == S.One
    assert sol.subs(n, 1) == S.One
\end{lstlisting}

\end{document}
